\scanpage{265}
\mspage{135}

longues de \deleted{\unsure{co}}suspension et de cosuspension dans
les catégories $\mathrm{HOT}_{W}^{\bullet}(S)$. (Plus exactement, ce sont
des théories relatives aux dérivateurs induits
$\mathrm{HOT}_{W,S}$, $\deleted{\unsure{\mathrm H}}\mathrm{HOT}_{W,S}^{\bullet}$ sur $S$, $I\unsure{-|\mapsto}\mathrm{HOT}_{W}(S\times I)$ resp.
$I\unsure{\mapsto|-}\mathrm{HOT}_{W}^{\bullet}(S\times I\ldots$) De plus, les carrés
cocartésiens et les carrés cartésiens relatifs
aux dérivateurs induits $\mathrm{HOT}_{W,S}$ sont
\uline{exacts}, i.e les suites \deleted{exactes} longues
de suspension et de cosuspension sont
exactes (axiomes d'exactitude à \deleted{gauche} \inserted{droite}
et à gauche $\mathrm{Der}\,6\,\mathrm{bis}$ et $\mathrm{Der}\,6$.)

2°) Le \deleted{pré}dérivateur $\mathrm{HOT}_{W}$ satisfait aussi
l'axiome $\mathrm{Der}\,1$ de \deleted{\unsure{somme directe}}\inserted{décomposition (finie $*$)}
% slanted
\margin{sans aucun axiome sur $W$ !}
(pour $I=\coprod_{\alpha}I_{\alpha}$),
\deleted{\unsure{Ainsi}} et l'axiome $\mathrm{Der}\,2$ de localisation.
De plus, $W$ est fortement saturé i.e. $f\in W\Longleftrightarrow\mathrm{hot}_{W}(f)$ iso.

3°) Concernant l'axiome $\mathrm{Der}\,3$ et ses variantes
(cf I, \unsure{pages|p.ges} 111--112, 125), on a \deleted{le résultat} ceci :

\deleted{\unsure{a)}} a) \deleted{La \unsure{signification} D-équivalences $W$}
% slanted
\margin{On n'a \deleted{pas} besoin ici \unsure{que} de $\mathrm{Loc}(1\,\mathrm{ter},\unsure{2},3)$ ($\mathrm{Loc}(1\,\mathrm{ter})$ est la saturation forte i.e. $\deleted{\unsure{a}}f\in W\Longleftrightarrow\mathrm{hot}_{W}(f)$ iso.}
Soit $f$ \inserted{$:S'\to S$} une flèche \unsure{de|dans} $\mathrm{Cat}$. Pour\textsuperscript{*}
\unsure{que} ce soit une $\mathrm{HOT}_{W}$-équivalence, i.e. pour
que le foncteur $f^{\unsure{*|!}}:\mathrm{HOT}_{W}(S)\to\mathrm{HOT}_{W}(S')$ induise une
\uline{équivalence} entre les sous-catégories pleines
formées des \deleted{\unsure{$f$}} objets \uline{constants} sur $S,S'$, il f. et
s. que $f$ soit dans $W$. \deleted{\unsure{On a}} L'énoncé correspondant dans les dérivateurs induits $\mathrm{HOT}_{W}(I)$
est \deleted{valable} également valable, pour $I\ne\varnothing$. Donc
à fortiori, $\mathrm{HOT}_{W}$ et $\mathrm{HOT}_{W}$ satisfont
% archive sheet number: 264
